\documentclass{beamer}

% Tema
\usetheme{Madrid}
\usecolortheme{default}

% Paquetes
\usepackage[utf8]{inputenc}
\usepackage{amsmath}
\usepackage{graphicx}

% Información
\title{Módulo 3: Construcción de Escenarios Macro-Financieros}
\author{MACROPOL}
\date{\today}

\begin{document}

%------------------------------------------------
\frame{\titlepage}

\begin{frame}{¿Qué haremos hoy y qué output producimos?}
\begin{itemize}
    \item Aprenderemos a distinguir tipos de escenarios en stress testing
    \item Construiremos una taxonomía aplicada al sistema previsional dominicano
    \item *Output de hoy:* Tabla de taxonomía de escenarios con variables relevantes para AFP/SIPEN
    \item *Conexión al pipeline:* Esta tabla es el esqueleto sobre el que calibraremos shocks en 3.2 y 3.3

\end{itemize}
    
\end{frame}
\begin{frame}{Frame Title}
El problema central: ¿qué le puede pasar al portafolio de una AFP?
Una AFP dominicana invierte en bonos soberanos, instrumentos del BC, y activos del exterior
Su portafolio es vulnerable a shocks de tasa de interés, tipo de cambio, inflación y crecimiento
Un stress test pregunta: *¿cuánto pierde ese portafolio si el entorno se deteriora?*
Para responder esa pregunta necesitamos primero *construir el entorno deteriorado* → eso es un escenario
Referencia: IMF Stress Testing Guide, Cap. 3; Farmer et al., Cap. 5

    
\end{frame}

\begin{frame}{Frame Title}
 ¿Qué es un escenario macroeconómico? (definición de trabajo)
Un escenario es un conjunto coherente de valores para variables macro-financieras bajo condiciones de estrés
No es una predicción: es una pregunta disciplinada del tipo *"¿qué pasa si...?"*
Tres requisitos: *plausibilidad* (puede ocurrir), *severidad* (es peor que lo normal), *coherencia interna* (las variables se mueven juntas con lógica económica)
Referencia: Siddique & Hasan, Cap. 2; IMF Stress Testing Guide
    
\end{frame}

\begin{frame}{Frame Title}
Taxonomía de escenarios: tres tipos que usaremos en este curso
Escenario histórico: reproduce un episodio real de estrés con sus magnitudes observadas (ej: crisis bancaria RD 2003, COVID-2020, crisis financiera global 2008)
Escenario hipotético: construido con lógica económica, no observado pero plausible (ej: caída del turismo + depreciación cambiaria + alza de tasas simultánea)
Escenario adverso severo: versión extrema, baja probabilidad pero alto impacto; exigido por reguladores (ej: SIPEN, FMI en evaluaciones FSAP)
*Los tres forman nuestro conjunto final de M3*
Referencia: Farmer et al., Cap. 5; IMF Stress Testing Guide, Cap. 4
    
\end{frame}

\begin{frame}{Frame Title}
Variables macro-financieras relevantes para el sistema dominicano
Tasa de interés: rendimiento de bonos soberanos DOP, tasa de política monetaria del BCRD
Tipo de cambio: DOP/USD (exposición de activos externos de las AFP)
Inflación: erosiona el valor real de las pensiones y afecta la tasa de descuento actuarial
Crecimiento del PIB: determina la masa salarial y por tanto las cotizaciones al sistema
Spread soberano: riesgo país que afecta el valor de mercado de los bonos en cartera
Referencia: SIPEN — Informes de Portafolio AFP;
    
\end{frame}
%------------------------------------------------
\end{document}